% Transcription — fonds Alexandre Grothendieck, shelfmark 86, batch 1 (pages 1-20).
% Established from the facsimile archives/batches/86/batch-01.pdf, cut from a
% copy of 86.pdf fetched through the project's own relay
% (https://grothendieck-relay.onrender.com/source/86.pdf), not uploaded by the
% user; PDF page k = archivists' page k, no cover sheet. Per the skill
% transcribe-grothendieck. The folder is thirty-two pages; this is the first
% of two batches. Pages 21-32 are batch 2, transcribed in parallel by another
% pass; its header was read at the end of this pass (see « tridron » below).
%
% Dating and title. The dating is the folder's own, copied verbatim from the
% catalogue; the group « Géométrie et topologie combinatoire » (69-89) holds
% it, and since the folder has a date of its own the group's range is not
% added. The inventory title « [Polyèdres réguliers] » is bracketed, i.e. the
% archivists'; it is kept, with its brackets, only in \foldertitle. His own
% words are on the three tan covers of this batch, and those alone become
% headings: p. 1 « Polyèdres réguliers I », p. 9 « Polyèdres réguliers II »
% with « géométrie d'incidence et ens. à opérateurs ... », p. 17 « Cours /
% Polyèdres réguliers III ». The archivists' title is presumably taken from
% them. The covers take no \page marker, by the rule for covers.
%
% Pages skipped: none apart from the three covers (1, 9, 17), which carry
% nothing but those inscriptions. Pages summarised: none (nothing typed,
% nothing by another hand). Pages transcribed from his hand: 2-8, 10-16,
% 18-20, all in blue felt-tip, with pencil in some drawings of pp. 5, 7, 8.
% Pages 6, 7 and 8 are mostly drawings of the icosahedron in various
% projections; they are described in French \note{}s, and only their text
% and formulas are set.
%
% His pagination: none on the leaves. On p. 10 he writes « cf. p. 1 », which
% presumably means p. 2 here (the first leaf of the run, where the ordered
% set Φ = S ⊔ A ⊔ F is introduced); it is recorded in a note, not resolved.
% The batch ends mid-sentence on p. 20 (« ... induisent »), continued at the
% top of p. 21 (batch 2 reads « des permutations paires sur l'ens des
% sommets »).
%
% What the batch is about. Three runs on the combinatorial icosahedron.
% I (pp. 2-8): p. 2 defines a « polyèdre combinatoire » as finite S, A, F with
% incidences R ⊂ S×A, R' ⊂ A×F, the ordered set Φ = S ⊔ A ⊔ F, R'' = R∘R',
% axioms 1) each edge has two vertices and two faces, 2) diamond property,
% 3)-3') the links Φ_{<f}, Φ_{>s} are connected (combinatorial polygons with
% q(f), p(s) ≥ 2 sides), 4) majorised sups exist, unpacked as a)-b), with a
% NB analysing when Sup(s,s'), Sup(s,a), Sup(a,a') exist, the dual 4'), and a
% Corollary that Φ → 𝔓(S) is an injective map of ordered sets. P. 3-4:
% every facet is the Sup of its vertices (Cor. Main 2), every face has at
% least three sides (Cor. 3), the Remark that the icosahedron is determined by
% S and A ⊂ 𝔓_2(S), F = triples all of whose pairs are edges, Aut = the
% permutations of S preserving A (the combinatorial distance). P. 5: a plan
% « Icosaèdre » (description, enumeration s = 12, a = 30, f = 20,
% 2a = ps = qf, p = 5, q = 3, incidence geometry / combinatorial polyhedra,
% combinatorial distances, perspectives by a pair of opposite vertices with
% symmetry group D_5 or Z_2 × D_5, transitivity on vertices, pairs (s<a),
% triples (s<a<f): the icosahedron is a « polyèdre combinatoire régulier »).
% Pp. 6-8: three explicit descriptions (A), (B), (C) of the icosahedron by
% its vertex set, adapted to an axis through two opposite vertices
% (S = {o} ⊔ Π ⊔ Π' ⊔ {o'}, two pentagons, symmetry Z/2 × D_5 of order
% 20 = 120/6), two opposite faces (S = T ⊔ T_1 ⊔ T'_1 ⊔ T', symmetry
% Z/2 × 𝔖_3 of order 12 = 120/10), and two opposite edges (S = σ ⊔ φ ⊔ σ' ⊔
% φ' ⊔ [σ×φ ⊔ σ'×φ']/..., symmetry Z_2 × (Z_2 × Z_2) of order 8 = 120/15),
% each with the list of edges by type and the count of each type (for (C)
% the counts as read do not obviously sum to 30; see the notes on p. 8).
% II (pp. 10-16, « Géométries d'incidence et ensembles à opérateurs »): for
% an incidence geometry Π = (S, A, F; R, R') of dimension 2, drapeaux and
% repères Rep(Π) = {(s,a,f) | s<a<f}; the three forgetful maps from Rep are
% of degree 2, giving fixed-point-free involutions σ0, σ1, σ2 with σ0σ2 =
% σ2σ0, hence an action of the group 𝔖_2 = <σ0, σ1, σ2 | σi² = 1, σ0σ2 =
% σ2σ0> (his letter; not the symmetric group); Theorem: the functor to
% 𝔖_2-sets is fully faithful (S, A, F recovered as quotients E/(σ1,σ2),
% E/(σ0,σ2), E/(σ0,σ1)); corollaries: Aut(Π) = Aut of the 𝔖_2-set Rep(Π);
% connectedness ⇔ transitivity; Aut acts freely; regularity (Aut transitive
% on Rep, card Aut = card Rep, the σi lifted to automorphisms);
% Isom(Π0, Π) ≃ Rep(Π); a cancelled definition of orientations R+ / R-; a
% Proposition (« Exer ») that p = ord(σ1σ2), q = ord(σ0σ1) are the vertex and
% face valencies, G_s ≃ D_p, G_f ≃ D_q, G_a ≃ Z/2 × Z/2.
% III (pp. 18-20, « Cours »): for the combinatorial icosahedron with its
% antipodism a, the set Ā of edge directions and a canonical permutation
% ρ : Ā → Ā (the edge of the pentagon around s opposite the second end t of
% a), ρ³ = id, five orbits called « trièdres distingués », orthogonal edge
% directions, the quotient T of cardinal 5; Proposition: orthogonal edges
% have no common vertex; P(s) ≃ T for each vertex; T is a torsor under the
% rotation group G_s^+ ≃ Z/5, which therefore acts by even permutations.
%
% Notation fixed once for the batch: his double-struck letters as \mathbf
% (\mathbf{Z}, \mathbf{D}_5); the powerset as \mathfrak{P}, \mathfrak{P}_2;
% the large gothic S of pp. 11-16 as \mathfrak{S}_2 (for his group on three
% generators, not the symmetric group; the \mathbf{G}_2 of p. 13 kept as
% written); the underlined a of p. 18 (antipodism) as \underline{a}; Π for
% both his pentagons (p. 6) and the incidence geometry (pp. 10 ff.), as he
% writes; Π_{s,u} for the vertices of pp. 8, 18, 20; Φ, Φ_{≥s}, Φ_{≤f} for
% the ordered set of facets and its links; Sup, Inf, Rep, Drap, Aut, Isom,
% tr as operator names.
%
% Relation to batch 2: the word read « trièdre » (pp. 19-20, « trièdre
% distingué ») is read « tridron » by the batch-2 pass; the ink (an
% unaccented five- or six-letter word) was re-cropped and left \uncertain at
% its first occurrence here. G_s^+ here is G_s in batch 2's summary.
%
% Readings to check first: the whole lower half of p. 2 (the NB 1°)-3°),
% the boxed (b), the oblique margin note) and the first line of p. 3; p. 4
% « Corollaire 4 »; p. 5 item 3 (a line of tangled insertions); p. 8, the
% whole (C) construction (the quotient by Z/2Z, σ(s,u), antipodism, the
% count 2); p. 11 the edge label on R' -> F; p. 12 the last sentence; p. 14
% the right-margin note to d)-d') and the replacement text for the cancelled
% Corollary; p. 15 the cancelled blocks; p. 18 the indices of the Π_{s,u}.
%
% Pass: Opus 5.5 (claude-opus-5-5), 2026-09-26 — first pass, unchecked against the pages by a human.

\documentclass[11pt,a4paper]{article}
% Le préambule commun à toutes les transcriptions du fonds Grothendieck.
%
% Il tient en une idée : **l'incertitude se compose, elle ne se résout pas.**
% Une transcription qui lisse un mot douteux a détruit la seule chose pour
% laquelle on la faisait — un lecteur doit pouvoir savoir, à chaque endroit,
% ce qui était sur le feuillet et ce qui a été ajouté. Les six macros qui
% suivent sont donc l'appareil critique, et non de la décoration.
%
% Les mêmes commandes sont reconnues par `scripts/render.mjs`, qui produit la
% vue de lecture du site. Une macro ajoutée ici doit l'être aussi là-bas,
% faute de quoi le rendu échouera bruyamment — ce qui est voulu : un rendu
% silencieusement faux est pire qu'un rendu absent.

\ProvidesPackage{grothendieck}[2026/08/08 Transcriptions du fonds Grothendieck]

\usepackage{amsmath,amssymb,amsthm}
\usepackage{mathrsfs}
\usepackage{stmaryrd}
% Les diagrammes commutatifs sont partout dans ce fonds : c'est la langue
% même de la géométrie algébrique de Grothendieck, et une transcription qui
% les rendrait en prose ne transcrirait rien.
\usepackage{tikz-cd}
% Français par défaut — c'est la langue des feuillets — mais l'anglais est
% chargé aussi : la traduction et le résumé sont des documents anglais, et la
% typographie française (espaces avant les deux-points) y serait une faute.
% Ces documents basculent par \selectlanguage{english} en tête de corps.
\usepackage[english,french]{babel}
\usepackage{fontspec}
\usepackage{geometry}
\geometry{a4paper,margin=25mm,marginparwidth=28mm}
\usepackage{marginnote}
\usepackage{xcolor}
% ulem, not soul. `soul`'s \st and \ul are built on a character-by-character
% scan that XeLaTeX's font machinery breaks: the document fails with an
% undefined control sequence rather than degrading. ulem does the same two jobs
% and is Unicode-engine safe. `normalem` stops it hijacking \emph, which the
% transcriptions use for Grothendieck's own emphasis.
\usepackage[normalem]{ulem}
\usepackage{hyperref}
\hypersetup{colorlinks=true,linkcolor=black,urlcolor=black,pdfborder={0 0 0}}

% --- Métadonnées du lot -----------------------------------------------------
% Reprises telles quelles par le script de rendu, qui les affiche en tête de
% la vue de lecture. Elles fixent ce que le fichier prétend couvrir : sans
% elles, une transcription flotte sans point d'attache dans un fonds de seize
% mille feuillets.

\newcommand{\folder}[1]{\gdef\grothFolder{#1}}
\newcommand{\batch}[1]{\gdef\grothBatch{#1}}
\newcommand{\leaves}[2]{\gdef\grothFirst{#1}\gdef\grothLast{#2}}
\newcommand{\foldertitle}[1]{\gdef\grothFoldertitle{#1}}
\gdef\grothFolder{?}\gdef\grothBatch{?}\gdef\grothFirst{?}\gdef\grothLast{?}\gdef\grothFoldertitle{}

% --- Le résumé --------------------------------------------------------------
%
% Il ouvre la lecture modernisée, sous le titre « Résumé », et il est composé
% en petit. Ce n'est pas un ornement : c'est le seul passage du document qui
% ne soit pas des mathématiques, et un lecteur doit pouvoir le reconnaître
% avant de l'avoir lu — puis le sauter s'il connaît déjà le sujet, ou s'y
% arrêter s'il ne le connaît pas. Le filet qui le suit marque la reprise du
% corps.

\newenvironment{resume}
  {\small\setlength{\parskip}{0.45em}}
  {\par\medskip\hrule\medskip}

% --- Mots-clés ---------------------------------------------------------------
%
% En anglais, en clôture du résumé de la lecture modernisée : le vocabulaire
% moderne sous lequel chercher ce que les feuillets construisent. Le manifeste
% du site (`npm run manifest`) les extrait pour étiqueter la cote entière —
% cette ligne est la source unique des tags, il n'y a pas de fichier à côté.
\newcommand{\keywords}[1]{\par\smallskip\noindent
  {\footnotesize\textsc{Keywords} --- \textit{#1}}\par}

% --- L'appareil critique ----------------------------------------------------

% Le numéro de feuillet, en marge. C'est l'ancre qui permet de régler un
% désaccord sur une formule en regardant *un* feuillet plutôt qu'une chemise
% de six cents. Le site s'en sert aussi pour tourner les pages du fac-similé
% quand on fait défiler la transcription.
\newcommand{\leaf}[1]{%
  \marginnote{\footnotesize\color{gray!70}\textsf{#1}}%
  \label{leaf:#1}\ignorespaces}

% Illisible : on ne devine pas. Un mot inventé qui se lit comme les autres est
% le pire résultat possible.
\newcommand{\ill}{\textcolor{red!60!black}{[\,\ldots\,]}}

% Lecture douteuse : le mot est donné, et signalé comme incertain.
\newcommand{\uncertain}[1]{\uline{#1}}

% Ajout de l'éditeur — une lettre manquante, un mot sous-entendu.
\newcommand{\add}[1]{\textcolor{blue!50!black}{[#1]}}

% Biffé par Grothendieck lui-même. Ce qu'il a rayé fait partie du document :
% une rature dit souvent où la pensée a changé de direction.
\newcommand{\struck}[1]{\sout{#1}}

% Note du transcripteur, sur l'état matériel du feuillet.
\newcommand{\note}[1]{\par\smallskip{\small\color{gray!60!black}\textit{[#1]}}\par\smallskip}

% Note marginale de Grothendieck, distincte du corps du texte.
\newcommand{\marginal}[1]{\par\smallskip\noindent\hspace{1em}%
  {\small\color{gray!60!black}#1}\par\smallskip}

% --- Titre ------------------------------------------------------------------

\AtBeginDocument{%
  \begin{center}
    {\large\bfseries Fonds Alexandre Grothendieck --- cote n\textsuperscript{o}~\grothFolder}\\[2pt]
    {\itshape\grothFoldertitle}\\[4pt]
    {\small feuillets \grothFirst--\grothLast\ (lot \grothBatch)}\\[2pt]
    {\footnotesize\color{gray!60!black}Transcription établie d'après le fac-similé numérisé
     par l'Université de Montpellier.\\
     Les lectures incertaines sont soulignées, les illisibles notées [\,\ldots\,],
     les ajouts entre crochets.}
  \end{center}
  \vspace{1em}}

\endinput


\folder{86}
\batch{1}
\pages{1}{20}
\dating{[à partir de 1977]}
\watermark{Édition de démonstration}
\foldertitle{[Polyèdres réguliers] : notes manuscrites (s.d.)}

\begin{document}

\section*{Polyèdres réguliers I}

\note{inscrit au crayon, en haut à droite de la chemise (page 1), qui ne
porte rien d'autre ; le « I » est souligné. La chemise ne reçoit pas de
numéro de page. C'est vraisemblablement de là que vient le titre entre
crochets de l'inventaire}

\subsection*{Polyèdres combinatoires}

\note{titre souligné, de sa main, en tête de la page 2}

\page{2}
$S$, $A$, $F$, $R$, $R'$ finis,
$R \subset S \times A$, $R' \subset A \times F$ (« relations d'incidence »)
$\Rightarrow$ $\Phi = S \amalg A \amalg F$ ens.\ ordonné ;
$R'' \subset S \times F$, $R'' = R \times_A R' = R \circ R'$.

\begin{enumerate}
  \item[1)] $\forall a \in A$, $\operatorname{card} R(a) = \operatorname{card}
  \uncertain{F(a)} = 2$ ;
  \item[2)] $\forall s \in S$, $\forall f \in F$ avec $s < f$,
  $\operatorname{card}\{a \in A \mid s < a < f\} = 2$ ;
  \item[3)] $\forall f \in F$, l'ens.\ ordonné $\Phi_{<f} = \{x \in \Phi \mid
  x < f\}$ (qui est un contour) est connexe, i.e.\ est un polygone
  combinatoire (à $q(f) \geq 2$ côtés) ;
  \item[3')] $\forall s \in S$, l'ens.\ ordonné $\Phi_{>s} = \{x \in \Phi \mid
  x > s\}$ (« id.\ ») est connexe, i.e.\ est un polygone combinatoire (à
  $p(s) \geq 2$ côtés) ;
  \note{« id. » rend les guillemets de répétition qu'il met sous
  « (qui est un contour) » de la ligne précédente}
  \item[4)] Dans $\Phi$ les sup majorés existent ($\Longleftrightarrow$ les
  inf minorés existent), i.e.\ $x, y \in \Phi$, $x \not\leq y$,
  $y \not\leq x$, $x, y < f \in F$ $\Rightarrow$ $\operatorname{Sup}(x,y)$
  existe,

  i.e.\ a) Toute arête est déterminée par l'ens.\ $\{s, s'\}$ de ses
  extrémités, et les faces qui majorent \uncertain{alors} majorent
  l'arête.

  b) Toute face $f$ est connue quand on connaît une arête $a$ de $f$ et un
  sommet de $f$ non incident, et \uncertain{bien} connaît deux sommets de
  $f$ \uncertain{non} incidents à une même arête.
\end{enumerate}
\note{le bloc 1)--3') est embrassé à gauche par une accolade, le long de
laquelle court une note verticale, non lue, qui se termine par
« connexité » (?)}

\textbf{NB} Indépendamment de 4), si $s, s' \in S$, $s \neq s'$,
$s, s' < a \in A$, \struck{alors si}
\begin{enumerate}
  \item[1°)] $\operatorname{Sup}(s,s')$ existe ssi $a = \operatorname{Sup}(s,s')$
  --- \struck{\ill{}} i.e.\ ssi \struck{\ill{}} a) ce soit seule arête qui
  soit incidente à $s, s'$ ; b) Toute face incidente à $s, s'$ l'est à $a$).
  \item[2°)] Si $s, s' \in S$, $s, s' < f \in F$, $s, s'$ non incidents à
  une \struck{\ill{}} même arête, alors $\operatorname{Sup}(s,s')$ existe
  $\Longleftrightarrow$ $f = \operatorname{Sup}(s,s')$ $\Longleftrightarrow$
  $f$ l'unique face qui contient $s, s'$.
  \item[3°)] Si $s \in S$, $a \in A$ non incidents, et $s, a < f \in F$, alors
  $\operatorname{Sup}(s,a)$ existe ssi $f = \operatorname{Sup}(s,a)$,
  \struck{\ill{}} i.e.\ ssi $f$ est l'unique face majorant $s$ et $a$.
  \item[3°)] Si $a \in A$, $a' \in A$, $a \neq a'$, $a, a' < f \in F$, alors
  $\operatorname{Sup}(a,a')$ existe ssi $f = \operatorname{Sup}(a,a')$ i.e.\
  ssi $f$ est l'unique face majorant $a, a'$.
\end{enumerate}
\note{deux alinéas numérotés « 3°) » sur la page ; sous 2°) une ligne
biffée commençant par « 3°) ». Les 1°)--3°) sont entourés d'un long trait}

(et assertions duales) $\Longleftrightarrow$ (b) Si $f, f' \in F$,
$f \neq f'$, alors l'ens.\ des facettes \uncertain{incidentes} à $f, f'$
est soit vide, soit réduit à un sommet, soit formé d'une arête et de ses
\uncertain{deux sommets} ; \uncertain{ou} $a \in A$ connu quand on connaît
l'ens.\ de ses sommets, i.e.\ $A \hookrightarrow \mathfrak{P}_2(S)$.
\note{cette dernière phrase est dans un encadré à droite, relié au bloc
1°)--3°) par une double flèche ; l'ordre de ses deux membres est
incertain}

4') (\uncertain{équiv.}\ à 4) a') Toute arête $a$ est connue quand on connaît
l'ens.\ $\{f, f'\}$ des deux faces incidentes, et tt sommet qui
\uncertain{est} \uncertain{incident} \struck{\ill{}} à celles-ci est
incident à l'arête.

b') Un sommet $s$ est connu quand on connaît une arête $a$ et une face $f$
\struck{incidentes} de $s$, non incidentes \ldots

\marginal{\uncertain{On a aussi} a') $A \hookrightarrow \mathfrak{P}_2(F)$,
b') Si $s, s' \in S$, $s \neq s'$, alors l'ens.\ des facettes \ill{} est
soit vide, soit \ill{} \ill{}}
\note{note écrite en oblique dans la marge inférieure gauche, dualisant
l'encadré (b) ; lecture d'ensemble douteuse}

Pour un \uncertain{polyèdre} combinatoire, \ldots

\textbf{Corollaire} 1°) Si : toute facette est \uncertain{associée} à l'ens.\
des sommets qu'elle majore, on trouve une application d'ens.\ ordonnés
\[
  \Phi \longrightarrow \mathfrak{P}(S)
\]
qui est injective, et induit un iso.\ de $\Phi$ sur
\note{la phrase se poursuit en haut de la page suivante}

\page{3}
\uncertain{l'ens.\ d'un ordonné} \struck{linéaire} une partie de $\mathfrak{P}(F)$
\uncertain{ordonnée} par inclusion --- partie stable par Sup de deux él.\
majorés \ldots\ Pour tt sommet \ill{} au moins trois arêtes et trois
faces.

\uncertain{Notons aussi} :

\textbf{Cor} Toute face est Sup des \uncertain{sommets} --- tt \ill{} est
Inf de ses \ill{}.

\uncertain{Connexité} \ill{} d'un \ill{} \uncertain{(par constr.\ de $R''$)},
par une suite de telles faces, on voit que
\[
  \operatorname*{Sup}_{a<f} a \;=\; \operatorname*{Sup}_{a<f}\,
  \operatorname*{Sup}_{s<a} s \;=\; \operatorname*{Sup}_{s<f} s \;=\; f
  \qquad \text{cqfd}
\]
\note{au-dessus du deuxième signe $=$ : « cor.\ 2 »}
(plus \uncertain{dual} \ldots)

\textbf{Remarque} Dans le cas de l'icosaèdre, \ill{} structure \ill{} \ill{}
connaît $S$ et $A \subset \mathfrak{P}_2(S)$ --- i.e.\ \ill{}
$F \subset \mathfrak{P}_3(S)$ : \uncertain{partie de} $C'$, comme l'ens.\ des
$f \in \mathfrak{P}_3(S)$ tels que $\mathfrak{P}_2(f) \subset A$, et le groupe
des automorph.\ de l'icosaèdre s'identifie au groupe des permutations $u$ de
$S$ telles que $u(A) = A$ (i.e.\ $u$ transforme arête en arête --- i.e.\
conserve la « distance combinatoire »).
\marginal{\ill{} \uncertain{des arêtes} \ill{} \uncertain{icosaèdre} \ill{}}
\note{note de quelques lignes écrite en diagonale dans la marge gauche, en
regard de la Remarque ; non lue}

\page{4}
un ens.\ de parties de $S$ ordonné par inclusion.

b) \struck{L'image} \uncertain{Cet} ens.\ de parties est stable par
intersection de deux parties \uncertain{majorées} :
$\operatorname{Im}(\Phi \to \mathfrak{P}(S))$.

\textbf{Dém.} b) résulte \uncertain{aussitôt} de l'existence des Inf de deux
él.\ minorés de $\Phi$. Pour a), \uncertain{et} le fait que
$\alpha(x) \leq \alpha(y) \Rightarrow x \leq y$ résulte \uncertain{dans}
\ldots

\textbf{Cor.\ Main 2} Toute face est le Sup des \struck{faces} sommets
qu'elle majore.

Pour sommet c'est trivial, pour arête c'est \uncertain{cela} résulte de
4 b) et de 4 a), pour face \uncertain{cela résulte de} \ldots

\textbf{Corollaire 3} Toute face a au moins trois côtés (i.e.\ a au moins
trois sommets).

En effet, s'il n'en avait que deux, les deux côtés auraient le même ens.\
de sommets, \uncertain{contrairement} à 4 a).

\textbf{NB} \uncertain{Si on se} \ill{} \uncertain{intersection} partie.
Les \uncertain{arêtes} de l'icosaèdre \ill{} comme les parties \ill{} les él.\
de $\mathfrak{P}_2(S)$, les \uncertain{faces} \ill{} les él.\ de
\uncertain{$\mathfrak{P}_3(S)$}.
\note{l'indice de ce dernier $\mathfrak{P}$ peut se lire 5 ; 3 est
ce qu'appelle la Remarque de la page~3}

\uncertain{Dualement}

\textbf{Corollaire \uncertain{4}} \uncertain{Tt sommet} est \struck{\ill{}}
\ill{} des faces qui \ill{}. \uncertain{L'involution} qui \ill{}
\ill{} des faces qui \uncertain{l'entourent}, définit une injection
$S \to \mathfrak{P}(F)$, induisant un \ill{}

\page{5}
\section*{Icosaèdre}

\begin{enumerate}
  \item[I] La description concrète et ses premières propriétés
\end{enumerate}

1. Description \uncertain{grossière} par un modèle concret (ou dessin)

2. Énumération. NB des sommets \quad $s = 12$\\
\phantom{2. Énumération. NB} des arêtes \quad $a = 30$ \qquad
$2a = ps = qf$\\
\phantom{2. Énumération. NB} des faces \quad $f = 20$ \qquad
$p = 5$, $q = 3$
\note{le premier chiffre de ce « 2. » est surchargé ; la ligne suivante
porte aussi « 2. »}

2. Description des propriétés polyédrales

\quad A) Description comme « géométrie d'incidence »

\quad B) \rule{3em}{0.4pt} « polyèdres combinatoires »

3. Distances combinatoires \uncertain{existence des} Sup et Inf
\uncertain{minorés}, \struck{propriétés des} \uncertain{des sommets},
\struck{\ill{}} \uncertain{caractérisation} de \uncertain{sommets}
\note{ligne très surchargée : insertions interlinéaires et ratures
enchevêtrées ; seuls « Distances combinatoires » et « Sup et Inf » sont
sûrs}

4. Perspectives

\quad A) Par sommet un couple sommet \uncertain{opposé} \ill{}
(groupe d'automorph.\ $\mathbf{D}_5$ ou $\mathbf{Z}_2 \times \mathbf{D}_5$)

\quad \textbf{Parenthèse} \uncertain{Groupe} des automorphismes des polygones
combinatoires \ldots

\quad Transitivité des groupes des automorphismes sur sommets, \ldots\ et
$=$ sur couples $(s,a)$ $(s<a)$ et $=$ sur triples $(s<a<f)$.

On dit que l'icosaèdre est un \emph{polyèdre combinatoire régulier}
(\textbf{NB} Il y a une infinité d'autres polyèdres combinatoires
réguliers --- mais parmi ceux-ci, \uncertain{seuls} \ill{} i.e.\
\uncertain{« réalisables »} comme des polyèdres convexes dans l'espace de
dimension 3 \ldots
\marginal{\uncertain{Interprétation} \uncertain{« antipodisme »} comme
automorphisme combinatoire}
\marginal{\ill{} \uncertain{polyèdre} \ill{}}
\note{deux notes en diagonale dans la marge gauche ; la seconde, au niveau
de A)--B), n'est pas lue. Au crayon, très pâle, un croquis d'icosaèdre
derrière le bas de la page}

\page{6}
\textbf{(A)}
\[
  S = \{o\} \amalg \Pi \amalg \Pi' \amalg \{o'\}
\]
($\Pi$, $\Pi'$ deux pentagones combinatoires, $\Pi \simeq \Pi'$,
$s \mapsto s'$)
\[
  A = \begin{cases}
    \{o, s\}, & s \in \Pi \\
    \{o', s'\}, & s' \in \Pi'
  \end{cases} \quad 10
  \qquad
  A(\Pi),\ A(\Pi') \quad 10
  \qquad
  \{s, t\},\ s \in \Pi,\ t \in \Pi' \quad 10
\]
$t$ incident à \uncertain{$s'$}, i.e.\ $\delta(t, s') = 1$
$\bigl(= \delta(t', s)\bigr)$.
\note{« $\delta(t',s)$ » est écrit sous « $\delta(t,s')$ », rattaché par un
signe d'égalité vertical}

Groupe de symétrie $\mathbf{Z}/2\mathbf{Z} \times \mathbf{D}_5$, d'ordre
$20 = 120/6$. \quad \textbf{NB} $6 = 12/2$.

\note{la page est occupée pour l'essentiel par des figures. En haut, (A$_1$) :
l'icosaèdre vu le long de l'axe $o o'$, décagone de sommets $s_0, \ldots,
s_4$ (autour de $o$, en trait plein) et $s'_0, \ldots, s'_4$ (autour de
$o'$, en tirets). Au milieu, (A$_2$), « ou » : deux pentagones emboîtés,
$s'_0, \ldots, s'_4$ à l'extérieur et $s_0, \ldots, s_4$ à l'intérieur
autour de $o$, les rayons issus des $s'_i$ marqués de flèches « $o'$ » (le
sommet $o'$ rejeté à l'infini). En bas, sous un trait horizontal, un petit
icosaèdre en perspective et un second, plus grand, dont les sommets sont
nommés $a, a', a''$, $b, b'$, $c$, $u, u'$, $w, w'$ (lettres d'identification parfois douteuses)}

\page{7}
\textbf{(B)}

Groupe de symétrie $\mathbf{Z}/2\mathbf{Z} \times \mathfrak{S}_3$, d'ordre
$12 = 120/10$. \quad \textbf{NB} $10 = 20/2$.

\note{page de figures. En haut, (B$_1$) : l'icosaèdre en perspective,
sommets nommés $u_1$, $u'$, $v'_1$, $w'_1$ (?), $v_1$, $w_1$, $u'_1$, \ldots,
vu le long d'un axe passant par les centres de deux faces opposées. En bas,
(B$_2$) : une construction en perspective de l'icosaèdre (sommets $a, a'$,
$b, b'$, $c, c'$, $u, v, w$), entourée de lignes de fuite, avec un triangle
hachuré ; à gauche, dans un cercle, une projection sur le plan d'une face,
sommets $u', v', w'$ au bord, $u_1, v_1, w_1$, $u'_1, v'_1, w'_1$ à
l'intérieur ; à droite, un icosaèdre inscrit dans un cube, à l'encre et au
crayon, et un petit croquis de cercles entrelacés}

\page{8}
\note{en haut à gauche, (B$_3$) : la même projection que la figure
circulaire de la page~7, au crayon, sommets $u'$, $v'$, $w'$
aux coins d'un grand triangle, $u_1, v_1, w_1$, $u'_1, v'_1, w'_1$, $u, v, w$
à l'intérieur ; au-dessus, un petit icosaèdre dans un cercle}
\[
  S = T \amalg T_1 \amalg T'_1 \amalg T'
\]
\note{le dernier terme est surchargé}
\[
  \begin{array}{ccccc}
    t & \boxed{T \xrightarrow{\ \sim\ } T'} & t' \\
    \wr & \wr \quad\ \wr & \wr \\
    t_1 & T_1 \simeq T'_1 & t'_1
  \end{array}
\]
($t \mapsto t'$, $t_1 \mapsto t'_1$) \struck{\ill{}} couple de 2 triangles en
corr.\ 1-1.

$A$ (arêtes) :
\[
  \begin{array}{lll}
    \mathfrak{P}_2(T), \quad \mathfrak{P}_2(T') & & 6 \\[2pt]
    \{t, u\},\ t \in T,\ u \in T_1,\ u \neq t_1 &
    \{t', u'\},\ t' \in T',\ u' \in T'_1,\ u' \neq t'_1 & 12 \\[2pt]
    \{t_1, u\},\ t_1 \in T_1,\ u \in T'_1,\ u \neq t'_1 & & 6 \\[2pt]
    \{t, t'_1\}, \quad \{t', t_1\} & & 6
  \end{array}
\]
\note{la première lettre de « $\{t_1, u\}$ » et de « $\{t, t'_1\}$ » est
surchargée ; les indices et accents de ces deux lignes sont incertains}

\textbf{(C)}
\[
  S = \sigma \amalg \varphi \amalg \sigma' \amalg \varphi' \amalg
  \bigl[\sigma \times \varphi \amalg \sigma' \times \varphi'\bigr]
  \big/ \uncertain{\mathbf{Z}/2\mathbf{Z}}
\]
\struck{$(s,u) \sim (t', v')$ si \ill{}}, $\sigma$ opérant par
\[
  \sigma(s,u) = \bigl((\varepsilon s)', (\varepsilon u)'\bigr),
\]
antipodisme $a(s,u) = (\varepsilon s, \varepsilon u)$ \ldots\
$(s,u) \in a \times b$ ou $(s,u) \in (a' \times b')$.
\note{le sens de $\sigma$ (ensemble à deux éléments et opération) et le
rôle de $a \times b$, $a' \times b'$ restent douteux : la ligne est
surchargée et l'appel qui les relie aux lignes au-dessus est une flèche}

$A$ :
\[
  \begin{array}{lll}
    \sigma, \sigma' & & \uncertain{2} \\
    \{s, u\},\ \{s', u'\} \quad (s \in \sigma,\ u \in \varphi) & & 8 \\
    \{\Pi_{s,u}, \Pi_{s,v}\},\ u \neq v & & 4 \\
    \{u, \Pi_{s,u}\},\ \{u', \Pi_{s',u'}\} & & 8 \\
    \{s, \Pi_{s,u}\},\ \{s', \Pi_{s',u'}\} & & 8 \\
    \{u, v'\},\ u \neq v & & 2
  \end{array}
\]
\note{à droite de la ligne $\{\Pi_{s,u}, \Pi_{s,v}\}$, un mot noirci ; les
deux premières lignes sont entourées et reliées par une flèche à la
troisième}

Automorph.\ $\mathbf{Z}_2 \times (\mathbf{Z}_2 \times \mathbf{Z}_2)$,
d'ordre $8 = 120/15$. \quad \textbf{NB} $15 = 30/2$.

\note{figures : (C$_1$), au crayon et à l'encre, l'icosaèdre vu le long
d'un axe passant par les milieux de deux arêtes opposées, sommets $u, v$,
$u', v'$, $s, t$, $\Pi_{s,u}$, $\Pi_{t,v}$, \ldots ; « ou » (C$_2$), la même
vue en projection, sommets extérieurs $s'$, $u'$, $v'$, $t'$, points
$\Pi_{t,u} = \Pi_{s',v'}$ (?), $\Pi_{t,v} = \Pi_{s',u'}$,
$\Pi_{s,u} = \Pi_{t',v'}$, $\Pi_{s,v} = \Pi_{t',u'}$ ; (C$_3$), un schéma
analogue plus petit, flèches vers $t'$ et $s'$}

\section*{Polyèdres réguliers II}

\note{inscrit au crayon sur la chemise (page 9), qui ne reçoit pas de
numéro de page : à gauche « Polyèdres réguliers », et dessous « II » ; à
droite « géométrie d'incidence et ens.\ à opérateurs \ldots »}

\page{10}
\subsection*{Géométries d'incidence et ensembles à opérateurs}

Soit $\Pi = (S, A, F ; R, R')$ une géométrie d'incidence de dim 2,
\struck{On appelle drapeau} (cf.\ p.~1 : ens.\ ordonné qu'il définit --- de
« dim.\ combinatoire » 2 ---). On appelle \emph{drapeau} de $\Pi$ une partie
tot.\ ordonnée non vide de $\Pi$ --- si son cardinal est $l$, on dit que le
drapeau est de \emph{longueur} $l-1$. Le maximum des longueurs des
drapeaux est 2 --- les drapeaux de longueur 2 sont appelés \emph{repères},
leur ens.\ se note $\operatorname{Rep}(\Pi)$
\[
  \operatorname{Rep}(\Pi) = \{(s, a, f) \in S \times A \times F \mid
  s < a < f\}
\]
\note{« cf.\ p.~1 » est de lui ; sa page 1 est vraisemblablement la page~2
de cette liasse, où l'ensemble ordonné $\Phi = S \amalg A \amalg F$ est
introduit}

Il y a des drapeaux des types suivants
\[
  \begin{array}{lll|l}
    0 & 1 & 2 & \text{drapeaux minimaux (de long.\ 0)} \\
    S & A & F & \text{ou « facettes »} \\[4pt]
    (0,1) & (1,2) & (0,2) & \text{drapeaux de long.\ 1} \\
    R \subset S \times A & R' \subset A \times F & R'' \subset S \times F & \\[4pt]
    (0,1,2) & & & \text{repères, ou drapeaux} \\
    \operatorname{Rep} \subset S \times A \times F & & & \text{maximaux}
  \end{array}
\]

Il y a des applications naturelles « d'oubli » (projections) entre ces
ens.\ de drapeaux, en particulier

\page{11}
\begin{tikzcd}[column sep=small, row sep=small]
  & \operatorname{Rep} \arrow[dl, "(\sigma_0)"'] \arrow[d, "(\sigma_1)"]
    \arrow[dr, "(\sigma_2)"] & \\
  \operatorname{Drap}_{12} = R' & R'' = \operatorname{Drap}_{02}
  & R = \operatorname{Drap}_{01}
\end{tikzcd}

avec $R' \subset A \times F$, $R'' \subset S \times F$, $R \subset S \times A$.
\note{à gauche du diagramme, un premier essai biffé ($\operatorname{Rep}$,
$R \subset S \times A$, $R' \subset$ \ldots)}

\note{ce qui suit est un ajout de sa main, encadré en haut à droite de la
page, entre crochets ; les « 2 » sont les degrés portés sur les flèches}
[En fait, on part de

\begin{tikzcd}[column sep=small, row sep=small]
  & R \arrow[dl, "2"'] \arrow[dr, "2"] & & R' \arrow[dl, "2"'] \arrow[dr] & \\
  S & & A & & F
\end{tikzcd}

et \ldots\ construit en dessus de ceux-ci $E = R \times_A R'$ d'où :

\begin{tikzcd}[column sep=small, row sep=small]
  & & E = \operatorname{Rep} \arrow[dl, "2"'] \arrow[dr, "2"] & & \\
  & R \arrow[dl, "2"'] \arrow[dr, "2"] & & R' \arrow[dl, "2"'] \arrow[dr, "2"] & \\
  S & & A & & F
\end{tikzcd}

(axiomes 1), 2))]
\note{sur la flèche $R' \to F$ du premier schéma, un signe surchargé, non
lu}

Ces trois applications sont de degré 2 --- d'où des involutions sans pt
fixe dans $\operatorname{Rep}$ notées $\sigma_0, \sigma_1, \sigma_2$ qui sont
caractérisées par le fait de conserver les fibres des applications
précédentes (\textbf{NB} $\sigma_i$ change la facette composante de dim $i$
d'un repère $r = (s, a, f)$ et \uncertain{conserve} les deux autres). On
voit que $\sigma_0$ et $\sigma_2$ \struck{\ill{}} commutent :
\[
  \sigma_0\sigma_2(s, a, f) = \sigma_2\sigma_0(s, a, f) = (s', a, f')
\]
où $s'$ est l'unique sommet $\neq s$ incident à $a$, $f'$ l'unique face
$\neq f$ incidente à $a$.

On voit que $\sigma_0, \sigma_1, \sigma_2$ et $\sigma_0\sigma_2$ opèrent sur
$\operatorname{Rep}$ sans pt fixe.

Soit
\[
  \mathfrak{S}_2 = \Bigl\{ \sigma_0, \sigma_1, \sigma_2 \Bigm|
  \begin{array}{l}
    \sigma_0^2 = \sigma_1^2 = \sigma_2^2 = 1 \\
    \sigma_0\sigma_2 = \sigma_2\sigma_0
  \end{array} \Bigr\}
\]
\note{la lettre du groupe est une grande lettre gothique, lue $\mathfrak{S}$
; l'indice 2 est sûr. Ce n'est pas le groupe symétrique $\mathfrak{S}_2$ :
c'est le groupe défini par la présentation qui suit}

\page{12}
Donc on a un foncteur
\[
  \begin{array}{c}
    \text{Géom.\ d'incidence} \\ \text{(avec iso)}
  \end{array}
  \longrightarrow
  \begin{array}{l}
    \mathfrak{S}_2\text{-ens.\ où } \sigma_0, \sigma_1, \sigma_2 \\
    \text{et } \sigma_0\sigma_2 \text{ opèrent sans pts} \\
    \text{fixes (avec iso)}
  \end{array}
\]
On va voir que \struck{le} \add{ce} foncteur est \emph{pleinement fidèle}
--- donc qu'on peut reconstruire $(S, A, F, R, R')$ à partir de
$E = \operatorname{Rep}$ en tant que $\mathfrak{S}_2$-ens.
\note{« Thm » est écrit au-dessus de « ce foncteur », comme titre de
l'énoncé}
Notons qu'on a
\begin{align*}
  E/(\sigma_1, \sigma_2) &\longrightarrow S = \Phi_0 \\
  E/(\sigma_0, \sigma_2) &\longrightarrow A = \Phi_1 \\
  E/(\sigma_0, \sigma_1) &\longrightarrow F = \Phi_2
\end{align*}
\struck{qui sont des applications surjectives.} Il est clair que
l'application médiane est un iso [$\sigma_0, \sigma_2$ engendrent un
s.-groupe de $\mathfrak{S}_2$ $\simeq \mathbf{Z}_2 \times \mathbf{Z}_2$, qui
opère librement sur $E$ \struck{\ill{}}, le quotient est
$\xrightarrow{\ \sim\ } A$]. Or les axiomes 3), et 3') (\uncertain{connexité})
expriment respectivement que la dernière et la première flèches sont
\struck{\ill{}} \emph{bijectives}. \uncertain{On peut donc exprimer} $S, A, F$ à partir de
$\uncertain{R}$.
\note{un mot souligné puis biffé entre « dernière et la » et « première »
; la dernière phrase est d'une lecture douteuse}

\page{13}
Quant à $R$, $R'$, ce sont respectivement les images de
$E = \operatorname{Rep}$ dans $S \times A$ et $A \times F$ \ldots

\textbf{Corollaire} \struck{Les automorph.\ de} Le groupe
$\operatorname{Aut}(\Pi)$ s'identifie à
$\operatorname{Aut}_{\mathfrak{S}_2\text{-ens}}(\operatorname{Rep}(\Pi))$.

\textbf{Corollaire} Conditions équivalentes
\begin{enumerate}
  \item[a)] Deux sommets de $\Pi$ \uncertain{sont} \add{toujours} reliés par
  une chaîne d'arêtes
  \item[a')] Deux faces de $\Pi$ sont \add{toujours} reliées par une chaîne
  d'arêtes
  \item[b)] $\Pi \simeq \Pi' \amalg \Pi'' \Longrightarrow \Pi' = \emptyset$
  ou $\Pi'' = \emptyset$
  \item[c)] \struck{\ill{}} $\mathfrak{S}_2$ opère transitivement dans $E$
  ($= \operatorname{Rep}(\Pi)$) (i.e.\ --- si $\Pi \neq \emptyset$ i.e.\
  $E \neq \emptyset$ --- $E$ est un espace homogène sous $G_2$).
\end{enumerate}
\note{« toujours » est deux fois une insertion interlinéaire (lue « tjs »).
Le « $G_2$ » de c) est en caractères gras, et diffère du $\mathfrak{S}_2$ de
la même ligne ; il est gardé tel quel}

\textbf{\struck{Corollaire}} \textbf{Déf} On dit alors que $\Pi$ est
\emph{connexe}.

\textbf{Corollaire} Si $\Pi$ est connexe, $\operatorname{Aut}\Pi$ opère
librement dans $\operatorname{Rep}(\Pi)$, i.e.\ si
$g \in \operatorname{Aut}(\Pi)$, $r \in \operatorname{Rep}(\Pi)$, alors
\[
  g \cdot r = r \Longrightarrow g = \mathrm{id}
\]
En effet, un automorphisme d'un espace homogène qui fixe un pt,
\struck{fixe \uncertain{chaque}} est l'identité. \struck{\ill{}}

\page{14}
\textbf{Corollaire} (si $\Pi$ fini :) $\operatorname{Card}\operatorname{Aut}(\Pi)
\mid \operatorname{Card}\operatorname{Rep}(\Pi)$

\textbf{Re-corollaire} Conditions équivalentes pour $\Pi$ connexe (et
\uncertain{soit} $r \in E = \operatorname{Rep}(\Pi)$ fixé)
\begin{enumerate}
  \item[a)] $\operatorname{Aut}(\Pi)$ est transitif sur
  $\operatorname{Rep}(\Pi)$
  \item[b)] Le stabilisateur de $r$ est un s.-groupe distingué de
  $\mathfrak{S}_2$
  \item[c)] (Si $\Pi$ est fini) $\operatorname{Card}\operatorname{Aut}(\Pi) =
  \operatorname{Card}\operatorname{Rep}(\Pi)$
  \item[d)] $\forall g \in \mathfrak{S}_2$, $\exists\, \bar g \in G$ tel que
  $g \cdot r = \bar g \cdot r$
  \item[d')] $\forall\, 0 \leq i \leq 2$, $\exists\, \bar\sigma_i \in G$ tel
  que $\sigma_i \cdot r = \bar\sigma_i \cdot r$
\end{enumerate}
\note{au-dessus de « $\operatorname{Aut}$ » en a), un « $G =$ » interlinéaire :
$G$ désigne $\operatorname{Aut}(\Pi)$ dans d) et d')}
\marginal{$\bar\sigma_i r \neq r$ \ldots\ $\bar\sigma_i$ \ill{} \ill{}}
\note{note serrée en marge droite de d)--d'), au plus petit module ; non lue
au-delà du premier signe}

\textbf{Déf} On dit alors que la géom.\ est \emph{régulière} (on parle
aussi de polyèdre combinatoire régulier \ldots)

\textbf{Cor.} Soient $(\Pi_0, r_0)$ polyèdre régulier avec repère, $\Pi$
polyèdre régulier isom.\ à $\Pi_0$. Alors
\[
  \operatorname{Isom}(\Pi_0, \Pi) \xrightarrow{\ \sim\ }
  \operatorname{Rep}(\Pi), \qquad \varphi \longmapsto \varphi(r_0)
\]
En particulier : Un couple $(\Pi_0, r_0)$, pour un type d'isomorphie donné
de $\Pi_0$, est défini à isom.\ unique près \ldots

\textbf{\struck{Corollaire}} \struck{$E/\sigma_0 \xrightarrow{\sim}
\operatorname{Drap}_{12}(\Pi)$, $E/\sigma_1 \xrightarrow{\sim}
\operatorname{Drap}_{0,2}(\Pi)$, $E/\sigma_2 \xrightarrow{\sim}
\operatorname{Drap}_{0,1}(\Pi)$} \struck{Soit} Avec les notations de d)
ci-dessus, $g \mapsto \bar g$ est un \uncertain{hom.}\ \ill{}
$\to G$, \ill{} \ill{}
\note{le « Corollaire » et ses trois isomorphismes sont barrés de croix ; la
phrase qui leur est substituée court dans la marge droite et se perd au
bas de la feuille}

\page{15}
\note{les deux blocs qui suivent, jusqu'à la Proposition, sont barrés de
longues diagonales ; ils restent lisibles et sont donnés en entier, la
biffure étant indiquée ici une fois pour toutes plutôt que ligne à ligne}
[Soit $\Pi = (S, \ldots)$ une géométrie d'incidence,
$R^+ \subset E = \operatorname{Rep}$, $R^- = E \setminus R^+$. Conditions
équivalentes
\begin{enumerate}
  \item[a)] $\sigma_i R^+ = R^-$ \quad ($i = 0, 1, 2$)
  \item[b)] $R^+ \xrightarrow{\ \sim\ } \Psi_i := \prod_{j \neq i} \Phi_j$
  \quad ($i = 0, 1, 2$) \quad où $\Phi_0 = S$, $\Phi_1 = A$, $\Phi_2 = F$
\end{enumerate}
On dit alors que $R^+$ est une \emph{orientation}, et les $r \in R^+$ sont
les \emph{repères directs} pour elle (ou appartenant à
l'\uncertain{orientation}), les $r \in R^-$ les \emph{repères inverses}]
\note{le $\prod$ de b) est lu tel quel}

Soit $\Pi$ polyèdre régulier connexe de type $(p,q)$.
\marginal{$p = \operatorname{ordre}(\sigma_1\sigma_2)_E$, $q =
\operatorname{ordre}(\sigma_0\sigma_1)_E$}
Groupe des automorph.\ d'un polyèdre régulier $\Pi$ qui fixent un sommet :
c'est le groupe diédral $\mathbf{D}_p$ (groupe des automorph.\ du polygone
\uncertain{plus précis} \ldots\ polygone) ; les $\bar\sigma$ engendrent $G$.
\note{ce second bloc est de plus encadré ; en marge gauche, un petit
croquis de polygone, les repères marqués $\sigma_0$, $\sigma_2$}

\marginal{(Exer)}
\textbf{Proposition} Soit $\Pi$ \struck{polygone}
polyèdre comb.\ régulier connexe, $E = \operatorname{Rep}(\Pi)$, $G =
\operatorname{Aut}(\Pi)$, $\Phi$ = ens.\ des facettes de $\Pi$ \ldots
\begin{enumerate}
  \item[a)] Soit $p = \operatorname{ordre}(\sigma_1\sigma_2)_E$,
  $q = \operatorname{ordre}(\sigma_0\sigma_1)_E$. Alors par tt sommet de
  $E\Pi$ \uncertain{passent} exactement $p$ arêtes, \uncertain{dans} tte
  face il y a exactement $q$ arêtes \\
  (\textbf{NB} $3 \leq p, q \leq +\infty$).
  \item[b)] Soit $s \in S(\Pi)$. Alors l'homom.\ du stabilisateur $G_s$
  de $s$ dans $G$, dans $\operatorname{Aut}(\Phi_{\geq s})$ [où
  $\Phi_{\geq s}$ est le polygone combinatoire à $p$ côtés
\end{enumerate}
\note{« $E\Pi$ » est écrit tel quel (le « $E$ » semble ajouté devant
« $\Pi$ »)}

\page{16}
formé des facettes incidentes à $s$] est un iso. En particulier,
$G_s \simeq \mathbf{D}_p$.
\begin{enumerate}
  \item[b')] Soit $f \in F(\Pi)$. Alors $G_f \to \operatorname{Aut}(\Phi_{\leq f})$
  (où $\Phi_{\leq f}$ est le polygone combinatoire à $q$ côtés des facettes
  incidentes à $f$) est un iso. En particulier, $G_f \simeq \mathbf{D}_q$.
  \item[c)] Soit $a \in A(\Pi)$. \struck{Alors} Soit $\sigma$ \add{(resp.\
  $\varphi$)} l'ens.\ des sommets (faces) incidents à $a$. --- donc $\sigma$,
  $\varphi$ sont des ens.\ de card.\ 2. Considérons
  \[
    G_a \longrightarrow \mathfrak{S}_\sigma \times \mathfrak{S}_\varphi
    = \mathbf{Z}/2\mathbf{Z} \times \mathbf{Z}/2\mathbf{Z}
  \]
  c'est un \emph{iso.}
\end{enumerate}
\note{le « (p) » au-dessus de « $\mathbf{D}_q$ » est une correction
interlinéaire douteuse. Le début de b) sur la page~15 s'achève
ainsi sur cette page}

\struck{\textbf{Exer} \textbf{Proposition} Soit $\Pi$ polyèdre régulier
\uncertain{comb.}\ connexe, \ill{} \ill{} $s \in S$, \ill{}}
\note{deux lignes barrées de croix en bas de page ; le reste de la feuille
est blanc}

\section*{Cours --- Polyèdres réguliers III}

\note{inscrit au crayon en haut à droite de la chemise (page 17), qui ne
reçoit pas de numéro de page : « Cours » souligné, puis « Polyèdres
réguliers », puis « III » (le chiffre souligné, lecture un peu douteuse)}

\page{18}
$\Pi = (S, A, F ; R, R', R'' ; E = \operatorname{Rep})$ \uncertain{icosaèdre}
combinatoire \\
$\underline{a}$ antipodisme, $\underline{a} \in G = \operatorname{Aut}(\Pi)$,
opère sans pts fixes sur $S, A, F$ (et a fortiori sur $R, R', R'', E$)
\[
  \bar A = A/\underline{a} \quad \text{ens.\ des paires « directions
  d'arêtes »} \ (\text{\uncertain{termes} \ill{} arêtes parallèles})
\]
$a \mapsto \bar a$ appl.\ canonique $A \to \bar A$.

On va définir une permutation canonique
\[
  \rho : \bar A \longrightarrow \bar A
\]
Pour ceci, pour $a \in A$, soit $s \in S$ incident à $a$ $= \{s, t\}$,
considérons les \struck{deux} arêtes $\{s, u\}$, $\{s, v\}$ \struck{dans}
\uncertain{de} $\Phi_{\geq s}$ adjacentes à $a$, considérons les deux arêtes
$\{s, \Pi_{s,u}\}$, $\{s, \Pi_{s,v}\}$ \uncertain{restantes} parmi les cinq de
$\Phi_{\geq s}$, donc $\{\Pi_{s,u}, \Pi_{s,v}\}$ est une arête de $\Pi$.
[On peut dire que les cinq sommets adjacents à $s$ forment un pentagone
\uncertain{pour} \uncertain{lequel} \uncertain{un sommet} distingué est le
deuxième sommet $t$ de $a$, et on prend le côté du pentagone opposé au
sommet $t$.] Si on avait pris l'autre sommet $t$, on aurait trouvé l'arête
$\{\Pi_{t,u}, \Pi_{t,v}\}$ antipodique de $\{\Pi_{s,u}, \Pi_{s,v}\}$, d'où un
unique élément $\rho_0(a) \in \bar A$. On trouve \uncertain{ainsi} une
application canonique
\[
  A \longrightarrow \bar A
\]
qui commute aux automorphismes de $\Pi$, et en particulier à
$\underline{a}$. Comme $\underline{a}$ opère trivialement sur $\bar A$,
$\rho_0$ se factorise en
\[
  \bar A \longrightarrow \bar A
\]
On trouve $\rho^2(\bar a) = \overline{\{u, v'\}} = \overline{\{v, u'\}}$
\note{en marge gauche, un croquis : l'arête $a = \{s, t\}$, les sommets
$u$, $v$, et les quatre sommets $\Pi_{s,u}$, $\Pi_{s,v}$, $\Pi_{t,u}$,
$\Pi_{t,v}$ aux coins d'un rectangle. Les indices des $\Pi$ sont serrés et
parfois surchargés ; ils sont donnés tels que l'argument les appelle
(cf.\ page~8, où apparaît la même notation $\Pi_{s,u}$)}

\page{19}
puis
\[
  \rho^3(\bar a) = \bar a
\]
Donc
\[
  \rho^3 = \mathrm{id}_{\bar A}
\]
En particulier, $\rho$ est une \emph{permutation}. Toutes ses orbites sont
de cardinal 3. Donc il y a cinq telles orbites.

Une orbite de $\rho$ s'appelle \emph{\uncertain{trièdre} distingué} (de directions
d'arêtes) associé à $\rho$. Deux directions d'arêtes $\alpha, \beta$ sont
dites \emph{orthogonales} si elles sont $\neq$ et dans le \uncertain{même}
\add{trièdre distingué} (i.e.\ dans le même trièdre distingué), ce qui
signifie aussi
\[
  \beta = \rho\alpha \quad \text{ou} \quad \alpha = \rho\beta
\]
Pour \uncertain{toute} \struck{\ill{}} direction d'arête, il y a deux
exactement deux directions d'arêtes orthogonales. [On notera que la
relation d'orthogonalité est symétrique par définition, mais la relation
$\beta = \rho\alpha$ \uncertain{entre $\alpha$ et $\beta$} ne l'est pas
évidemment.]
\note{le mot lu « trièdre » s'écrit sans accent visible, en cinq ou six
lettres ; la transcription du batch 2 le lit « tridron ». « deux exactement deux » est sur la page ; l'insertion au-dessus de
« relation » est de lecture douteuse. La ligne sous « dites orthogonales »,
« i.e.\ dans le même trièdre distingué », est interlinéaire}

L'ens.\ des trièdres distingués est un ens.\ $T$ \struck{\ill{}} quotient de
cardinal 5 de $\bar A$, noté $T$. \uncertain{\struck{On a}} appl.\ canonique
$\bar A \xrightarrow{\ \mathrm{tri}\ } T$. On peut le considérer comme un
ens.\ \uncertain{quotient} de $A$, formé de 5 parties de cardinal 6,
chacune d'elles stable par $\underline{a}$.
\textbf{Interprétation} concrète en termes de coloriage,
\marginal{\uncertain{terminologie} : \uncertain{arêtes} \uncertain{orthogonales} \ill{}}
\note{la ligne s'arrête sur une virgule, au bas de la feuille ; la page
suivante commence par une Proposition. Le « quotient » de l'avant-dernière
phrase est à cheval sur une note marginale écrite en oblique}

\page{20}
\textbf{Proposition} \struck{\ill{}} Deux arêtes orthogonales n'ont pas de
sommet en commun.

En effet, les \uncertain{quatre} arêtes orthogonales \uncertain{à}
$\{s, t\}$ \uncertain{sont} $\{\Pi_{s,u}, \Pi_{s,v}\}$, et $\{\Pi_{t,u}, \Pi_{t,v}\}$ d'une part,
$\{u, v'\}$ et $\{v, u'\}$ d'autre part, elles ne rencontrent pas $\{s, t\}$.

\textbf{Corollaire} Pour tt $s \in S$, si $P(s)$ est l'ens.\ des arêtes
incidentes à $s$, \struck{des $t \in S$ adjacents à $s$,} l'application
$a \mapsto \operatorname{tr}(\bar a)$ induit une bijection
\[
  P(s) \longrightarrow T
\]
En effet, c'est une application injective entre deux ens.\ de cardinal 5.

\textbf{Corollaire} Soit $G_s^+$ le groupe des rotations de $\Pi$ fixant
$s$ (i.e.\ les automorph.\ fixant $s$, et induisant sur $\Phi_{\geq s}$ une
rotation \ldots) Alors \struck{\ill{}} $T$ est un torseur sous $G_s^+$ ---
donc si on fixe un gén.\ $u_s$ de $G_s^+$, $u_s$ est une permutation
circulaire sur $T$.

Comme $T$ est de card.\ impair, on en conclut que $G_s^+$ opère par
\emph{permutations paires}.

\textbf{NB} Plus généralement, on voit que les automorph.\ d'un polygone
combinatoire de \uncertain{long.}\ \struck{impaire est dans}
$\equiv 1\ (4)$ (p.ex.\ d'un pentagone) induisent
\note{la phrase se poursuit au-delà de ce batch (page 21)}

\end{document}
